\documentclass[11pt,a4paper]{article}
\usepackage{fontspec}
\setmainfont{Latin Modern Roman}
\usepackage[margin=23mm]{geometry}
\usepackage{amsmath,amssymb,amsthm,mathtools,mathrsfs}
\usepackage{longtable,booktabs}
\usepackage[unicode,hidelinks]{hyperref}
\providecommand{\C}{\mathbb C}
\providecommand{\R}{\mathbb R}
\newtheorem{theorem}{Theorem}[section]
\newtheorem{lemma}[theorem]{Lemma}
\newtheorem{proposition}[theorem]{Proposition}
\newtheorem{corollary}[theorem]{Corollary}
\newtheorem{definition}[theorem]{Definition}
\newtheorem{remark}[theorem]{Remark}
\allowdisplaybreaks
\emergencystretch=3em
\makeatletter
\renewcommand\@pnumwidth{2.5em}
\renewcommand\@tocrmarg{3.5em}
\makeatother
\title{The Original Kernel Volume Bound}
\author{Split-Zero research programme}
\date{14 September 2026}
\begin{document}
\maketitle
\paragraph{Portable source edition.}
Every TeX input used by this entry is expanded below. Historical file
and directory names in the preserved text are source locators; the
exact origin paths and hashes are recorded in \texttt{11\_SOURCE\_MAP.json}.
The other numbered proof files in this flat handoff supply the
accompanying complete calculations.\par

% Complete finite estimate on the original m=1 packet and its actual kernel.
% Original source and quotient providers: IFB1--11 and GMB1--16.
\section{A finite bound for the original kernel volume from its actual dimension}
\label{sec:OKV}

\paragraph{The original packet, sources, and exact coordinate map.}
Throughout this calculation the original multiplicity is $m=1$ and
\[
 k=4l+1,\quad l\ge1,\quad q=(k+1)^2\ge36,\quad c=k/2,
 \qquad 0<\delta<1/2,\quad\gamma>2.
 \tag{OKV1}
\]
Retain every original root and its two coordinates:
\[
\begin{split}
 \chi(S)&=\prod_{a,b=0}^k
  [S-c-(2a-k)\delta-i(2b-k)\gamma],\\
 \omega_{ab}&=(2b-k)\gamma-i(2a-k)\delta,\\
 P(y)&=i^{-q}\chi(c+iy)=\prod_{a,b=0}^k(y-\omega_{ab}).
\end{split}
 \tag{OKV2}
\]
The identity follows factor by factor from
$iy-(2a-k)\delta-i(2b-k)\gamma=i(y-\omega_{ab})$;
it retains the full phase $i^q$. The algebra map
\[
 \Phi:\mathbb C[y]/(P)\longrightarrow E_\chi:=\mathbb C[S]/(\chi),
 \qquad [r(y)]\longmapsto[r((S-c)/i)]
 \tag{OKV3}
\]
is well defined because $P((S-c)/i)=i^{-q}\chi(S)$.
Its inverse is $[f(S)]\mapsto[f(c+iy)]$, as verified by substitution.
Both maps preserve degrees and are inverse on each polynomial ring
before taking the quotients. In particular division by $P$ is the
original remainder operation through this exact coordinate map.

Use only the two original source orders $s\in\{1,k\}$ and write
\[
 b_s=s/2,\qquad M_s=(2\pi)^{s/2},\qquad
 dm_{0,s}(y)=\frac{M_s2^{b_s-2}}{\pi\Gamma(b_s)}
     \left|\Gamma\left(\frac{b_s}{2}+\frac{iy}{2}\right)\right|^2dy.
 \tag{OKV4}
\]
The exact source calculation IFB2--5 supplies mass $M_s$, the transform
$\int e^{ty}dm_{0,s}=M_s(\cos t)^{-b_s}$ for $|t|<\pi/2$, and
the original monic orthogonal polynomials and norms
\[
\begin{gathered}
 p_0=1,\quad p_1=y,\qquad
 p_{d+1}=yp_d-d(b_s+d-1)p_{d-1},\\
 h_d^{(s)}=\int|p_d|^2dm_{0,s}=M_s d!(b_s)_d,\qquad
 \varphi_d(y)=p_d(y)/\sqrt{h_d^{(s)}},\qquad
 u_d(S)=\varphi_d((S-c)/i).
\end{gathered}
 \tag{OKV5}
\]
All these integrals are on the same real line, with the original
source mass. For example the stated orthogonality and norms also
follow directly from its transform and the exact generating function
\[
 \sum_{d\ge0}p_d(y)\frac{z^d}{d!}
       =(1+z^2)^{-b_s/2}e^{y\arctan z}.
\]
Indeed the source integral of the product at sufficiently small real
$z,w$ is $M_s(1-zw)^{-b_s}$. The cosine addition identity and the
displayed transform give this equality; the source exponential bound
justifies each finite coefficient derivative under the integral.
The coefficients of $z^dw^e$ then give zero for $d\ne e$ and
$M_s d!(b_s)_d$ for $d=e$. The initial coefficients are $1,y$,
and differentiating the generating function gives precisely the
recurrence in OKV5. Thus the source norms used below retain their
original constants and all finite-degree factors.

Let $C_s:\mathbb C^q\to E_\chi$ have columns $[u_0],\ldots,[u_{q-1}]$
in the original frame $(1,S,\ldots,S^{q-1})$. Its exact leading
coefficients give
\[
 \det C_s=i^{-q(q-1)/2}\prod_{d=0}^{q-1}(h_d^{(s)})^{-1/2}\ne0.
\]
Retain the original GMB vectors and matrices, including the last row:
\[
 v_j=C_s^{-1}[u_{q+j}],\qquad
 K_j=I_q+\sum_{a=0}^{j-1}v_av_a^*,\qquad 0\le j\le q+1.
 \tag{OKV6}
\]
The vector $v_q$ has original polynomial degree $2q$.

\paragraph{Exact remainder coefficients in all original roots.}
Index the full list in OKV2 by $\omega_1,\ldots,\omega_q$, with
each pair $(a,b)$ occurring exactly once. Define the elementary and
complete homogeneous polynomials by their finite sums
\[
 e_h(\omega)=\sum_{1\le a_1<\cdots<a_h\le q}
                      \omega_{a_1}\cdots\omega_{a_h},\qquad
 h_t(\omega)=\sum_{n_1+\cdots+n_q=t}\prod_{a=1}^q\omega_a^{n_a},
 \tag{OKV7}
\]
with $e_0=h_0=1$, $e_h=0$ outside $0\le h\le q$, and
nonnegative integer $n_a$ in the second sum. The notation $h_t(\omega)$
in this paragraph is a polynomial in the roots; the original squared
norm continues to be $h_d^{(s)}$ of OKV5. Multiplication of formal
power series gives the exact identity
\[
 \left(\sum_{h=0}^q(-1)^he_h(\omega)z^h\right)
 \left(\sum_{t\ge0}h_t(\omega)z^t\right)=1,
\]
because each geometric series $(1-\omega_a z)^{-1}$ multiplies its
own factor $1-\omega_a z$. Consequently, for $D=q+j$ with $0\le j\le q$,
the quotient and remainder of $y^D$ upon division by $P$ are exactly
\[
\begin{split}
 Q_D(y)&=\sum_{t=0}^jh_t(\omega)y^{j-t},\qquad
 \operatorname{rem}_P(y^D)=y^D-P(y)Q_D(y),\\
 [y^a]\operatorname{rem}_P(y^D)
   &=-\sum_{t=0}^j h_t(\omega)(-1)^{D-t-a}e_{D-t-a}(\omega),
       \qquad 0\le a<q.
\end{split}
 \tag{OKV8}
\]
To verify the quotient assertion, the power-series identity makes
the coefficients of $y^D,y^{D-1},\ldots,y^q$ in $P Q_D$ equal
to those of $y^D$. The remaining coefficients are exactly the last
line of OKV8. Their degree is less than $q$, so uniqueness of monic
polynomial division proves both assertions. For $0\le D<q$ the
remainder is $y^D$ itself.

The actual maximal root modulus is
\[
 R=k\sqrt{\delta^2+\gamma^2},\qquad
 D_*:=\max\{3,\sqrt{\delta^2+\gamma^2}\},\qquad T_*:=D_*q.
 \tag{OKV9}
\]
Every $|\omega_{ab}|\le R$, and equality holds at each corner,
so this expression for $R$ is exact. In particular $R\le T_*$.
Counting the summands of OKV7 proves
\[
 |e_h(\omega)|\le\binom qh R^h,\qquad
 |h_t(\omega)|\le\binom{q+t-1}{t}R^t.
\]
The second count follows by placing $q-1$ separators among $t+q-1$
positions to encode $(n_1,\ldots,n_q)$. For fixed $j$, the binomial
$\binom{q+t-1}{t}$ is increasing in $0\le t\le j$, because its
successive ratio is $(q+t)/(t+1)\ge1$. In each nonzero summand
of OKV8 the total root exponent is $(D-t-a)+t=D-a$.
As $t$ varies the indices $D-t-a$ are distinct, so the sum of
their $\binom q{D-t-a}$ is at most $2^q$. It follows that
\[
 \left|[y^a]\operatorname{rem}_P(y^{q+j})\right|
 \le2^q\binom{q+j-1}{j}R^{q+j-a}
 \le2^{3q-1}R^{q+j-a},\qquad 0\le a<q,\quad0\le j\le q.
 \tag{OKV10}
\]
For the last inequality use
$\binom{q+j-1}{j}\le2^{q+j-1}\le2^{2q-1}$, by the binomial
expansion of $(1+1)^{q+j-1}$. These estimates retain the exact
coefficient formula OKV8, including every signed cancellation.

\paragraph{Original Gamma norms of the low monomials.}
The same recurrence and norms in OKV5 give, without changing the measure,
\[
 y\varphi_a=\sqrt{(a+1)(b_s+a)}\,\varphi_{a+1}
          +\sqrt{a(b_s+a-1)}\,\varphi_{a-1},
 \quad \varphi_{-1}=0.
 \tag{OKV11}
\]
For a polynomial of degree at most $q-2$, its orthonormal coefficient
vector is therefore sent by multiplication by $y$ to the sum of
one upward and one downward weighted shift, with output degree at most
$q-1$. Each shift has operator norm at most
$\sqrt{q(b_s+q)}$: the squared norm of its coefficient vector is
bounded by the largest squared weight times the input squared norm.
Since $s\le k<q$ gives $b_s\le q/2$, the triangle inequality yields
\[
 \|yf\|_{L^2(m_{0,s})}
       \le2\sqrt{q(b_s+q)}\,\|f\|_{L^2(m_{0,s})}
       \le3q\,\|f\|_{L^2(m_{0,s})}
       \le T_*\|f\|_{L^2(m_{0,s})}.
\]
Here $2\sqrt{3/2}=\sqrt6<3$. Starting from
$\|1\|_{L^2(m_{0,s})}=\sqrt{M_s}$ and applying this inequality
to $1,y,\ldots,y^{a-1}$ proves
\[
 \boxed{\|y^a\|_{L^2(m_{0,s})}\le\sqrt{M_s}\,T_*^a,
                      \qquad0\le a<q.}
 \tag{OKV12}
\]
All intermediate polynomial degrees satisfy the stated operator domain.

For any polynomial $f=\sum_a f_a y^a$, write
$\mathfrak c_{T_*}(f)=\sum_a|f_a|T_*^a$. This is a finite
coefficient sum, not a change of source or quotient norm.
The exact original recurrence gives
\[
 \mathfrak c_{T_*}(p_{d+1})
 \le T_*\mathfrak c_{T_*}(p_d)
       +d(b_s+d-1)\mathfrak c_{T_*}(p_{d-1}).
\]
For $1\le d<2q$ its actual positive coefficient is bounded by
$d(b_s+d-1)\le2q(b_s+2q)\le5q^2\le T_*^2$.
The initial values are $1,T_*$. Induction now proves
\[
 \boxed{\mathfrak c_{T_*}(p_d)\le(2T_*)^d,
                            \qquad0\le d\le2q.}
 \tag{OKV13}
\]
Indeed substituting the two preceding bounds in the recurrence gives
$T_*(2T_*)^d+T_*^2(2T_*)^{d-1}
 =\tfrac34(2T_*)^{d+1}\le(2T_*)^{d+1}$.

Combining all $q$ coefficients of OKV10 with OKV12, by the triangle
inequality in the original $L^2$ space, gives
\[
 \|\operatorname{rem}_P(y^D)\|_{L^2(m_{0,s})}
           \le\sqrt{M_s}\,q2^{3q}T_*^D\quad(0\le D\le2q).
\]
For $D\ge q$, each summand is bounded by
$\sqrt{M_s}2^{3q-1}R^{D-a}T_*^a
 \le\sqrt{M_s}2^{3q-1}T_*^D$ and there are $q$ summands.
For $D<q$, the assertion follows from OKV12 and $q2^{3q}\ge1$.
Linearity of the exact remainder map, followed by OKV13, proves
\[
 \boxed{\|\operatorname{rem}_P(p_d)\|_{L^2(m_{0,s})}
         \le\sqrt{M_s}\,q2^{3q}(2T_*)^d,
                           \qquad0\le d\le2q.}
 \tag{OKV14}
\]

\paragraph{The original remainder vectors and the full matrix norm.}
The polynomial $\operatorname{rem}_P(p_d)$ has degree less than $q$.
Its expansion in the orthonormal basis
$\varphi_0,\ldots,\varphi_{q-1}$ has squared coefficient norm equal
to its squared original $L^2(m_{0,s})$ norm, by expanding the finite
inner product and using OKV5. The exact quotient map OKV3 and
definition OKV6 therefore give, for $d=q+j$, $0\le j\le q$,
\[
 \|v_j\|_2^2
   =\frac{\|\operatorname{rem}_P(p_{q+j})\|_{L^2(m_{0,s})}^2}
                 {M_s(q+j)!(b_s)_{q+j}}
   \le\frac{q^22^{6q}(2T_*)^{2(q+j)}}{(q+j)!(b_s)_{q+j}}.
 \tag{OKV15}
\]
This equality is the actual map of the original quotient coordinates
to the remainder polynomial. The full source factor $M_s$ occurs in
both the numerator bound and its original norm denominator before
the displayed cancellation. No estimate on the conditioning of
$C_s$ or a companion matrix is used.

Since $b_s\ge1/2$, for each integer $d\ge1$ the first factor of
$(b_s)_d$ is at least $1/2$ and every subsequent factor is at least
one. Thus $d!(b_s)_d\ge1/2$. Also $T_*\ge1$ and $q+j\le2q$.
Consequently
\[
 \|v_j\|_2^2\le2q^22^{6q}(2D_*q)^{4q},\qquad0\le j\le q,
\]
and the original complete matrix obeys
\[
 \boxed{1\le\|K_{q+1}\|_2\le
   1+2(q+1)q^22^{6q}(2D_*q)^{4q}.}
 \tag{OKV16}
\]
To check its norm bound, $\|v_jv_j^*\|_2=\|v_j\|_2^2$
follows directly from Cauchy--Schwarz, with equality on a nonzero
$v_j$; it is zero for a zero vector. Add all $q+1$ summands of
OKV6 and the norm of $I_q$. Positivity also proves its lower bound.

Here is an explicit bound of the requested scale, with constants
depending only on the original two root coordinates:
\[
 \boxed{\log\|K_{q+1}\|_2
       \le C(\delta,\gamma)\,q\log(q+1),\qquad
 C(\delta,\gamma)=12+4\log(2D_*).}
 \tag{OKV17}
\]
For a direct verification, put
$U=2(q+1)q^22^{6q}(2D_*q)^{4q}\ge1$ and use
$\log(1+U)\le\log(2U)$. This upper bound is exactly
\[
 \log4+\log(q+1)+2\log q+6q\log2
       +4q\log(2D_*)+4q\log q.
\]
Write $L=\log(q+1)$. Since $q\ge36$, one has
$\log4\le L$, $\log q\le L$, $6\log2\le2L$, and $L\ge1$.
The last inequality follows already from
$\log4=\int_1^4dt/t>1$, by bounding the integrand below by
$1/2$ on $[1,2]$ and $1/4$ on $[2,4]$.
The preceding display is therefore at most
$4L+6qL+4\log(2D_*)qL$, which is bounded by the right side
of OKV17. The completely finite and sharper bound OKV16 remains
available before this logarithmic estimate.

\paragraph{The four original kernel Grams with their full coefficient frame.}
Retain the original surjection $\Lambda:E_\chi\to B_{\rm obs}$,
the actual subspace $\mathcal K_\Lambda=\ker\Lambda$, and its original
inclusion $I:\mathcal K_\Lambda\hookrightarrow E_\chi$. Put
$r_K=\dim\mathcal K_\Lambda$; no separate formula or upper bound
for this actual dimension is needed here. Use its fixed original
coefficient frame and write
\[
 T_s=C_s^{-1}I,\qquad
 G_{q+j-1}^{(s)}=C_s^{-*}K_j^{-1}C_s^{-1},\qquad
 H_j^{(s)}=I^*G_{q+j-1}^{(s)}I=T_s^*K_j^{-1}T_s,
             \qquad0\le j\le q+1.
 \tag{OKV18}
\]
For completeness, the quotient metric in this identity follows from
the original orthonormal source-coordinate map
$A_j=[I_q,v_0,\ldots,v_{j-1}]$.
It obeys $A_jA_j^*=K_j$. The minimum norm solution of $A_jz=x$
is $z=A_j^*K_j^{-1}x$: it has the required image and is orthogonal
to $\ker A_j$. Every other solution differs by such a kernel vector,
so the Pythagorean identity gives minimum squared norm $x^*K_j^{-1}x$.
The actual remainder is $C_sx$. Substitution proves the original
quotient Gram and then its full restriction in OKV18.

Let $\lambda_s=\|K_{q+1}\|_2$. The original update chain gives
$I_q=K_0\preceq K_j\preceq K_{q+1}\preceq\lambda_s I_q$.
For positive definite matrices, inversion reverses the order:
if $A\preceq B$, conjugate by $A^{-1/2}$ to obtain a matrix
with eigenvalues at least one, invert those eigenvalues and conjugate
back. Thus the full original kernel forms satisfy
\[
 \lambda_s^{-1}H_0^{(s)}\preceq H_j^{(s)}\preceq H_0^{(s)},
 \qquad H_0^{(s)}=T_s^*T_s,
 \qquad H_{j+1}^{(s)}\preceq H_j^{(s)}.
 \tag{OKV19}
\]
For $r_K>0$, the matrix $T_s$ is injective, so $H_0^{(s)}$ is
positive definite. Conjugation in OKV19 by its inverse square root
gives eigenvalues in $[\lambda_s^{-1},1]$. The determinant identity
for that exact congruence is
\[
 \det\bigl[(H_0^{(s)})^{-1/2}H_j^{(s)}(H_0^{(s)})^{-1/2}\bigr]
       =\frac{\det H_j^{(s)}}{\det(T_s^*T_s)}.
\]
Therefore
\[
 -r_K\log\lambda_s
 \le\log\det H_j^{(s)}-\log\det(T_s^*T_s)\le0.
 \tag{OKV20}
\]
In particular the determinant and conditioning of the actual
$T_s^*T_s$ are retained in this identity; no new kernel basis is
substituted for its original frame.

Define the original four-endpoint kernel volume, with its literal
degrees $q-1,q,2q-1,2q$, by
\[
 F_K^{(s)}=
 \log\det H_0^{(s)}+\log\det H_1^{(s)}
       -\log\det H_q^{(s)}-\log\det H_{q+1}^{(s)}.
 \tag{OKV21}
\]
Monotonicity in OKV19 makes each of
$\log\det H_0^{(s)}-\log\det H_q^{(s)}$ and
$\log\det H_1^{(s)}-\log\det H_{q+1}^{(s)}$ nonnegative.
Each is at most $r_K\log\lambda_s$ by OKV20. Consequently the
actual kernel volume satisfies the complete finite estimates
\[
 \boxed{\begin{split}
 0\le F_K^{(s)}
 &\le2r_K\log\|K_{q+1}\|_2\\
 &\le2r_K\log\left[1+2(q+1)q^22^{6q}(2D_*q)^{4q}\right]\\
 &\le2r_K C(\delta,\gamma)q\log(q+1),
                    \qquad s\in\{1,k\}.
 \end{split}}
 \tag{OKV22}
\]
For $r_K=0$, every kernel determinant has its original empty value
one, and $F_K^{(s)}=0$. Hence OKV22 includes this case.
The four copies of $\log\det(T_s^*T_s)$ in OKV21 have coefficients
$+1,+1,-1,-1$ and cancel exactly only after their presence in
OKV20 has been recorded.

\paragraph{The original boundary and arithmetic transitions.}
At each original degree set
\[
 Q_j^{(s)}=
   \left(\Lambda (G_{q+j-1}^{(s)})^{-1}\Lambda^*\right)^{-1},
 \qquad
 S_j^{(s)}=(G_{q+j-1}^{(s)})^{-1}\Lambda^*Q_j^{(s)}.
 \tag{OKV23}
\]
Their zero dimensional cases have their empty meaning. Surjectivity
of $\Lambda$ and positivity give the displayed inverses. Direct
multiplication proves $\Lambda S_j^{(s)}=I$ and
$I^*G_{q+j-1}^{(s)}S_j^{(s)}=0$.
For any fixed original section $S_*$ of $\Lambda$, the map $[I,S_*]$
is bijective: its inverse records $\Lambda z$ and the unique kernel
coordinate of $z-S_*\Lambda z$. The minimum section differs from
$S_*$ by a kernel column. The two frames therefore differ by a
block triangular matrix with identity diagonal and determinant one.
Taking the determinant of the full metric in the orthogonal frame
$[I,S_j^{(s)}]$ gives
\[
 \det H_j^{(s)}\det Q_j^{(s)}
       =|\det[I,S_*]|^2\det G_{q+j-1}^{(s)}.
 \tag{OKV24}
\]
This proves the exact relation with all kernel and boundary directions
retained. The same frame factor occurs at all four endpoints.
Their signed sum is consequently
\[
 F_K^{(s)}+F_B^{(s)}=\mathcal B_k^{(s)},\qquad
 F_B^{(s)}=\log\det Q_0^{(s)}+\log\det Q_1^{(s)}
               -\log\det Q_q^{(s)}-\log\det Q_{q+1}^{(s)}.
 \tag{OKV25}
\]
The original matrices $G_{q+j-1}^{(s)}$ decrease with $j$, so their
inverse forms increase. Compression by $\Lambda$ and inversion then
show that $Q_j^{(s)}$ decreases as well. Thus $F_B^{(s)}\ge0$,
and the finite companion bound is
\[
 \max\{0,\mathcal B_k^{(s)}-2r_K C(\delta,\gamma)q\log(q+1)\}
       \le F_B^{(s)}\le\mathcal B_k^{(s)}.
 \tag{OKV26}
\]
Here the finer middle bound in OKV22 can be used in place of its
last line, with the same proof.

Finally the original source identifications are
$\mathcal B_k^{(1)}=\mathcal B_k^\sigma$ and
$\mathcal B_k^{(k)}=\mathcal B_k^0$. Their already proved full
source transition is $\mathcal H_k=\mathcal B_k^\sigma-\mathcal B_k^0$,
and the original arithmetic transition is
$\mathcal B_k^{\rm ar}=\mathcal B_k^\sigma+\delta_k^\sigma$,
with its original finite signed allowance for $\delta_k^\sigma$.
Substitution of the proved exact splitting OKV25 gives, with every
source and arithmetic term still present,
\[
 \boxed{\mathcal B_k^{\rm ar}
   =F_K^{(k)}+F_B^{(k)}+\mathcal H_k+\delta_k^\sigma
   =F_K^{(1)}+F_B^{(1)}+\delta_k^\sigma.}
 \tag{OKV27}
\]
Thus OKV22 controls the original Gamma kernel summands in this
same arithmetic identity. It keeps the separate source transition
and arithmetic scalar with their established signs and bounds.
The result is finite for the actual $r_K=\dim\ker\Lambda$ and
uses no additional asserted dimension estimate or replacement source.

\end{document}
